\documentclass[10pt]{article}
\usepackage[utf8]{inputenc}
\usepackage[T1]{fontenc}
\usepackage{amsmath}
\usepackage{amsfonts}
\usepackage{amssymb}
\usepackage[version=4]{mhchem}
\usepackage{stmaryrd}
\usepackage{bbold}

\begin{document}
$L 14$

L14: Proof of Quadratic Reciprocity Law (Theorem 4.9)\\
Lemma 4.10: Let $p$ be an odd prime number and let $a \in \mathbb{Z}$ with $p+a$ and $a$ odd. Let

$$
N:=\sum_{j=1}^{(p-1) / 2}\left[\frac{j a}{p}\right] .
$$

Then,

$$
\left(\frac{a}{p}\right)=(-1)^{N} .
$$

Here, $[\cdot]$ denotes the floor function.

\begin{itemize}
  \item Lemma 4.10, similar to Euler's criterion and Gauss' (emma) is of more theoretical interest rather than computational.\\
Example: Use Lemma 4.10 to compute $\left(\frac{7}{11}\right)$.\\
We have
\end{itemize}

$$
\begin{aligned}
N=\sum_{j=1}^{5}\left[\frac{j-7}{11}\right] & =\left[\frac{7}{11}\right]+\left[\frac{14}{11}\right]+\left[\frac{21}{11}\right]+\left[\frac{28}{11}\right] \\
& +\left[\frac{35}{11}\right] \\
& =0+1+1+2+3=7 .
\end{aligned}
$$

Thus $\left(\frac{2}{11}\right)=(-1)^{2}=-1$.

Proof (of Lemma 4.10): Let $r_{1}, \ldots, r_{n}$ be the least non-negative residues of the integers $a, 2 a, \cdots,\left(\frac{p-1}{2}\right)$ a that are greater than $P / 2$, and let $S_{1}, S_{2}, \ldots, S_{m}$ be the least non-negative residues of these integers that are less than $P / 2$. Then for $j=1,2, \ldots, \frac{P-1}{2}$, we have that\\
(1) $j a=p\left[\frac{j a}{p}\right]+$ remainder depending on $j$.\\
Each of the $r_{1}, r_{2}, \ldots, r_{n}, s_{1}, s_{2}, \ldots, s_{m}$ appears exactly once as a remainder.\\
By adding the $\frac{p-1}{2}$ equations in (1), we olotain

$$
\sum_{j=1}^{(p-1) / 2} j a=\sum_{j=1}^{(p-1) / 2^{2}} P\left[\frac{j a}{P}\right]+\sum_{j=1}^{n} r_{j}+\sum_{j=1}^{m} S_{j}
$$

From the proof of Ganss' Lemma, we\\
know that the integers $p-r_{1}, p-r_{2}, \ldots$, $p-r_{n}, S_{1}, S_{2}, \ldots, S_{m}$ are precisely the integers from I to $\frac{p-1}{2}$ inclusive in Some order. Thus $\sum_{j}^{(p-1) / 2} j=\sum_{j}^{n}\left(p-r_{j}\right)+\sum_{j}^{m} s_{j}=p n-\sum_{j}^{n} r_{j}+\sum_{j}^{m} s_{j}$\\
$j=1$

$$
j=1 \quad j=1 \quad j=1 \quad j=1
$$

(3).

Subtracting $_{(p-1) / 2}$ (3) from (2) we obtain

$$
\sum_{j=1}^{(p-1) / 2} j a-\sum_{j=1}^{(p-1) / 2} j=\sum_{j=1}^{(p-1) / 2} p\left[\frac{j a}{p}\right]-p n+2 \sum_{j=1}^{n} r_{j}
$$

Reducing both sides of (4) modulo 2, we have

$$
O \equiv \sum_{j=1}^{(p-1) / 2}\left[\frac{j a}{p}\right]-n(\bmod 2)
$$

and so $\quad(p-1) / 2$

$$
n \equiv \sum_{j=1}\left[\frac{j a}{p}\right](\bmod 2) .
$$

Thus $n \equiv N(\bmod 2)$, and so $\left(\frac{a}{p}\right)=(-1)^{n}$\\
(by Gauss Lemma) implies that $\left(\frac{a}{p}\right)=(-1)^{N}$.

Proof (of Theorem 4.9-Quadratic reciprocity):

We will prove that $\left(\frac{p}{q}\right)\left(\frac{q}{p}\right)=(-1)^{\frac{p-1}{2} \cdot \frac{q-1}{2}}$.\\
Without loss of generality we assume that $p>q$. Consider Figure 4.1 on $p g 122$ of the textbook. We will count, in two different ways, the number of lattice points in the rectangle $O A B C$, including those points on $A B$ and $B C$, but not those $O A$ and $O C$. This number is clearly $\frac{(p-1)}{2} \cdot \frac{(q-1)}{2}$. We will count a second way now. We have the following facts:\\
(a) ON has slope $q / p$. Since $p$ and $q$ are distinct prime numbers, there are no lattice points on $O N$ except for the endpoints $O$ and $N$.\\
(b) The $y$-coordinate of $M$ is

$$
\frac{q}{p} \cdot \frac{p-1}{2}=\frac{q}{2}-\frac{1}{2} \frac{q}{p} .
$$

(c) The $y$ co-ordinate of $M$ in (b) lies between two consecutive integers $\frac{q-1}{2}$ and $\frac{q+1}{2}$, since $\frac{q-1}{2}=\frac{q}{2}-\frac{1}{2}<\frac{q}{2}-\frac{7}{2} \frac{q}{p}$ (since $p>q)$

$$
\begin{aligned}
& <q / 2 \\
& <\frac{q+1}{2} .
\end{aligned}
$$

Thus there are no lattice points in the triangle $M P B$, except possibly on $P B$.

Thus the number of lattice points in the rectangle $O A B C$ including those points on $A B$ and $B C$, but not those points on $O A$ and $O C$, is the number of lattice points in the triangle $O P C$, not including those points on $O C_{3}$ plus the number of lattice points in the triangle $O A M$, not inclucling those points on $O A$. The number of lattice points\\
in the triangle OPC, not including those (6) points on $O C$, is

$$
N_{1}=\sum_{j=1}^{(q-1) / 2}\left[\frac{j p}{q}\right] .
$$

The number of lattice points in the triangle\\
$O A M$, not including those points on $O A$, is

$$
N_{2}=\sum_{j=1}^{(p-1) / 2}\left[\frac{j q}{p}\right] .
$$

Thus $(q-1) / 2, \quad(p-1) / 2$


\begin{align*}
N_{1}+N_{2} & =\sum_{j=1}\left[\frac{j p}{q}\right]+\sum_{j=1}\left[\frac{j q}{p}\right] \\
& =\frac{(p-1)}{2} \cdot \frac{(q-1)}{2} \text { by our first } \tag{5}
\end{align*}


Then by Lemma 4.10 we have that $\left(\frac{p}{q}\right)\left(\frac{q}{p}\right)=(-1)^{N_{1}}(-1)^{N_{2}}$

$$
\begin{aligned}
& =(-1)^{N_{1}+N_{2}} \\
& =(-1)^{\frac{p-1}{2} \cdot \frac{q-1}{2}} \text { by (S). }
\end{aligned}
$$

Example: Characterise the primes $P$\\
(in terms of congruence conditions) for which 11 is a quadratic residue modulo $p$.\\
We want all the primes such that $\left(\frac{11}{p}\right)=1$.\\
Using quadratic reciprocity we have that

$$
\left(\frac{11}{p}\right)= \begin{cases}\left(\frac{p}{11}\right) & \text { if } p \equiv 1(\bmod 4) \\ -\left(\frac{p}{11}\right) & \text { if } p \equiv 3(\bmod 4)\end{cases}
$$

Qnadratic vesidues $\bmod 11: 1,3,4,5,9(\bmod 11)$\\
Quadratic non-residues $\bmod 11: 2,6,7,8,10(\bmod 11)$\\
Then use the Chinese Remainder Theorem to solve each of the 10 pairs of congruences.

$$
\begin{aligned}
& \left\{\begin{array}{l}
p \equiv 1(\bmod 4) \\
p \equiv 1(\bmod 11)
\end{array} \quad \begin{array}{l}
\text { Exercise } \\
\Rightarrow
\end{array} p \equiv 1(\bmod 44)\right. \\
& \left\{\begin{array}{l}
p \equiv 1(\bmod 4) \\
p \equiv 3(\bmod 11)
\end{array} \Rightarrow p \equiv 25(\bmod 44)\right. \\
& \left\{\begin{array}{l}
p \equiv 1(\bmod 4) \\
p \equiv 4(\bmod 11)
\end{array} \Rightarrow p \equiv 37(\bmod 44)\right.
\end{aligned}
$$

$$
\begin{aligned}
& \left\{\begin{array}{l}
p \equiv 1(\bmod 4) \\
p \equiv 5(\bmod 11)
\end{array} \Rightarrow p \equiv 5(\bmod 44)\right. \\
& \left\{\begin{array}{l}
p \equiv 1(\bmod 4) \\
p \equiv 9(\bmod 11)
\end{array} \Rightarrow p \equiv 9(\bmod 44)\right. \\
& \left\{\begin{array}{l}
p \equiv 3(\bmod 4) \\
p \equiv 2(\bmod 11)
\end{array} \Rightarrow p \equiv 35(\bmod 44)\right. \\
& \left\{\begin{array}{l}
p \equiv 3(\bmod 4) \\
p \equiv 6(\bmod 11)
\end{array} \Rightarrow p \equiv 39(\bmod 44)\right. \\
& \left\{\begin{array}{l}
p \equiv 3(\bmod 4) \\
p \equiv 7(\bmod 11)
\end{array} \Rightarrow p \equiv 7(\bmod 44)\right. \\
& \left\{\begin{array}{l}
p \equiv 3(\bmod 4) \\
p \equiv 8(\bmod 11)
\end{array} \Rightarrow p \equiv 19(\bmod 44)\right. \\
& \left\{\begin{array}{l}
p \equiv 3(\bmod 4) \\
p \equiv 10(\bmod 4)
\end{array} \Rightarrow p \equiv 43(\bmod 44)\right.
\end{aligned}
$$

For primes $p$ we have $\left(\frac{11}{p}\right)=1$ if and only if $P \equiv 1,5,7,9,19,25,35,37,39,43(\bmod 44)$.

Example (Exercise): Characterise the primes $P$ for which $S$ is a quadratic residne modulo $P$.


\end{document}